% Part: computability
% Chapter: recursive-functions
% Section: general-recursive-functions

\documentclass[../../../include/open-logic-section]{subfiles}

\begin{document}

\olfileid{cmp}{rec}{gen}
\olsection{সাধারণ পুনরাবৃত্ত অপেক্ষক}

সর্বত্র-সংজ্ঞায়িত অপেক্ষকের একটি সেট পাওয়ার আরেকটি উপায় আছে। কোনো
সর্বত্র-সংজ্ঞায়িত অপেক্ষক $f(x,\vec z)$-কে \emph{নিয়মিত} বলা যাক,
যদি স্বাভাবিক সংখ্যার প্রত্যেক অনুক্রম $\vec z$-এর জন্য এমন একটি $x$
থাকে, যাতে $f(x,\vec z) = 0$। অন্যভাবে বললে, নিয়মিত অপেক্ষকগুলি ঠিক
সেই অপেক্ষক, যাদের উপর সীমাহীন অনুসন্ধান প্রয়োগ করে শেষে একটি
সর্বত্র-সংজ্ঞায়িত অপেক্ষক পাওয়া যায়। রক্ষণশীলভাবে সীমাহীন অনুসন্ধানকে
নিয়মিত অপেক্ষকের মধ্যেই সীমাবদ্ধ রাখা যায়:

\begin{defn}
\ollabel{defn:general-recursive}
স্বাভাবিক সংখ্যা থেকে স্বাভাবিক সংখ্যায় বিভিন্ন স্থানসংখ্যার অপেক্ষকের
যে ক্ষুদ্রতম সেট শূন্য, উত্তরসূরি ও অভিক্ষেপকে ধারণ করে এবং মিশ্রণ, আদিম
পুনরাবৃত্তি ও \emph{নিয়মিত} অপেক্ষকের উপর প্রয়োগ করা সীমাহীন অনুসন্ধানের
অধীনে বদ্ধ, তাকে \emph{সাধারণ পুনরাবৃত্ত অপেক্ষকসমূহের} সেট বলা হয়।
\end{defn}

স্পষ্টতই প্রত্যেক সাধারণ পুনরাবৃত্ত অপেক্ষক সর্বত্র-সংজ্ঞায়িত।
\olref{defn:general-recursive} ও \olref[par]{defn:recursive-fn}-এর
পার্থক্য হলো, পরের সংজ্ঞাটিতে মাঝের ধাপগুলিতে আংশিক পুনরাবৃত্ত অপেক্ষক
ব্যবহারের অনুমতি আছে; একমাত্র শর্ত, শেষে যে অপেক্ষকটি পাওয়া যাবে সেটি
সর্বত্র-সংজ্ঞায়িত হতে হবে। তাই “সাধারণ” শব্দটি একটি ঐতিহাসিক অবশেষ এবং
অযথার্থ নাম; আপাতদৃষ্টিতে \olref{defn:general-recursive},
\olref[par]{defn:recursive-fn}-এর চেয়ে \emph{কম} সাধারণ। কিন্তু সৌভাগ্যবশত
পার্থক্যটি অলীক; সংজ্ঞা দুটি আলাদা হলেও সাধারণ পুনরাবৃত্ত অপেক্ষকের সেট
ও পুনরাবৃত্ত অপেক্ষকের সেট অভিন্ন।

\end{document}
